\section{Caso de Aplicación IV: Comparación de Calidad de Proveedores de Grafeno mediante PCA y Clasificación con Kernel Ridge en Espectros FTIR}
% ==========================================

\subsection{Contexto y Fundamento del Control de Calidad}
En el ámbito industrial de los nanomateriales, la reproducibilidad lote a lote es un desafío crítico. Específicamente, en la síntesis de Óxido de Grafeno Reducido (rGO), el nivel de reducción química determina directamente la conductividad eléctrica y las propiedades mecánicas finales del material. El análisis directo de espectros FTIR mediante métodos convencionales puede verse dificultado por el ruido instrumental y la deriva de la línea base. 

Este caso de aplicación aborda la automatización del control de calidad y la clasificación de proveedores mediante técnicas estadísticas multivariantes: Análisis de Componentes Principales (PCA) mediante Descomposición en Valores Singulares (SVD) para reducción de dimensionalidad, y un clasificador Kernel Ridge con kernel RBF bajo un esquema Multiclase One-vs-Rest (OvR) para validar la separabilidad de los lotes.

\subsection{Modelo de Simulación Físico-Química de Reducción}
Para simular de manera realista la variabilidad en el control de calidad, la biblioteca genera espectros ficticios basados en un espectro real de óxido de grafeno ($y_{\text{base}}$) medido en el vector de números de onda $x$. Se introducen variaciones no colineares en la presencia de grupos funcionales de oxígeno para modelar de forma preciso la cinética química de reducción, donde los grupos hidroxilo (O-H), carbonilo (C=O) y alcoxilo (C-O) se reducen a distintas tasas físicas:

\begin{equation}
    y_{\text{sim}} = y_{\text{base}} \times s + \psi_{\text{quim}} + \text{drift}_{\text{lineabase}} + \epsilon_{\text{ruido}} + \delta_{\text{shift}}
\end{equation}

donde:
\begin{itemize}
    \item $s$ es un factor de escala aleatorio.
    \item $\psi_{\text{quim}}$ es la firma química específica de la clase del proveedor, construida sobre perfiles gaussianos para O-H ($3400\text{ cm}^{-1}$), C=O ($1720\text{ cm}^{-1}$) y C-O ($1050\text{ cm}^{-1}$):
    \begin{equation}
        P_{3400} = e^{-\left(\frac{x - 3400}{250}\right)^2}, \quad P_{1720} = e^{-\left(\frac{x - 1720}{50}\right)^2}, \quad P_{1050} = e^{-\left(\frac{x - 1050}{60}\right)^2}
    \end{equation}
    \item $\text{drift}_{\text{lineabase}} = c_1 x_{\text{norm}} + c_2 x_{\text{norm}}^2$ representa la deriva de la línea base debido a dispersión Mie residual.
    \item $\epsilon_{\text{ruido}} \sim \mathcal{N}(0, \sigma^2)$ es el ruido gaussiano instrumental.
    \item $\delta_{\text{shift}}$ es un desplazamiento vertical aleatorio del espectro completo.
\end{itemize}

Los tres proveedores se modelan con las siguientes características físicas y químicas (que eliminan la colinealidad lineal y reflejan el grado de reducción de grafeno):
\begin{enumerate}
    \item \textbf{Proveedor A (Alta Calidad - rGO altamente reducido)}: Presenta un grado óptimo de reducción. Los picos de grupos funcionales de oxígeno son prácticamente nulos: $\psi_{\text{quim}} = 0.05 P_{3400} + 0.05 P_{1720} + 0.05 P_{1050}$. El ruido es muy bajo ($\sigma = 0.005$) y la deriva de línea base es insignificante ($c_1, c_2 \in [-0.02, 0.02]$).
    \item \textbf{Proveedor B (Calidad Media - rGO parcialmente reducido)}: Muestra una reducción incompleta. Conserva grupos carbonilo y alcoxilo, pero tiene baja presencia de hidroxilos: $\psi_{\text{quim}} = 0.1 P_{3400} + 2.0 P_{1720} + 1.0 P_{1050}$. El nivel de ruido es moderado ($\sigma = 0.03$) y la deriva de línea base es intermedia ($c_1, c_2 \in [-0.15, 0.15]$).
    \item \textbf{Proveedor C (Baja Calidad - GO sin reducir)}: Representa material de baja calidad técnica (prácticamente óxido de grafeno puro sin reducir). Tiene una alta abundancia de hidroxilos y grupos oxigenados: $\psi_{\text{quim}} = 4.0 P_{3400} + 0.5 P_{1720} + 3.0 P_{1050}$. El nivel de ruido es severo ($\sigma = 0.09$) y presenta una deriva extrema de línea base ($c_1, c_2 \in [-0.40, 0.40]$).
\end{enumerate}

\subsection{Análisis de Componentes Principales (PCA) mediante SVD}
Para el PCA, primero se normalizan vectorialmente los espectros crudos $y_{\text{sim}}$ de forma que su norma $L_2$ sea la unidad. La matriz de datos espectrales resultantes $X$ ($N_{\text{muestras}} \times L$) se centra restando la media de cada variable, obteniéndose $X_c$.

La descomposición de valores singulares (SVD) de $X_c$ es:
\begin{equation}
    X_c = U \Sigma V^T
\end{equation}
donde las filas de $V^T$ representan las direcciones de los componentes principales (vectores de carga o \textit{loadings}). Las puntuaciones espectrales (\textit{scores}) se calculan proyectando los datos centrados:
\begin{equation}
    T = X_c V
\end{equation}
La varianza asociada al componente principal $i$-ésimo se obtiene a partir del valor singular correspondiente $\lambda_i = \Sigma_{ii}$:
\begin{equation}
    \sigma^2_i = \frac{\lambda_i^2}{N_{\text{muestras}} - 1}
\end{equation}
El porcentaje de varianza explicada se define como $\text{EVR}_i = \lambda_i^2 / \sum_j \lambda_j^2$.

\subsection{Resultados y Clasificación}
Al aplicar PCA sobre el dataset simulado de 180 espectros (60 por proveedor), el primer componente principal (PC1) captura el \textbf{55.36\%} de la varianza total, y el segundo (PC2) el \textbf{40.01\%} (explicando en conjunto el \textbf{95.37\%} de la varianza total). Como se observa en la figura \ref{fig:pca_suppliers}, el espacio bidimensional PC1-PC2 revela una clara segregación en tres grupos bien definidos correspondientes a cada proveedor:
\begin{itemize}
    \item El \textbf{Proveedor A} (verde) se concentra en una región estrecha y densa, debido a su bajísima variabilidad y ruido.
    \item El \textbf{Proveedor B} (naranja) se desplaza a lo largo del eje PC2 por la presencia selectiva de picos a $1720\text{ cm}^{-1}$.
    \item El \textbf{Proveedor C} (rojo) muestra una gran dispersión espacial debido al elevado nivel de ruido aleatorio e interferencias de línea base de gran magnitud, pero claramente separado por su dominante banda O-H a $3400\text{ cm}^{-1}$.
\end{itemize}

\begin{figure}[h!]
    \centering
    \includegraphics[width=0.7\textwidth]{pca_scores_ftir_quality.png}
    \caption{Espacio de puntuaciones de los componentes principales PC1 y PC2 para la clasificación de proveedores de grafeno (Caso IV).}
    \label{fig:pca_suppliers}
\end{figure}

Para validar computacionalmente esta separación, se entrenó y evaluó un clasificador Kernel Ridge con kernel RBF ($\gamma=0.1, \alpha=0.1$) en un esquema One-vs-Rest (OvR). Utilizando una validación cruzada de 5 pliegues (5-fold Cross-Validation), el clasificador alcanzó una precisión del \textbf{100.00\%}. Esta robustez demuestra que, a pesar del severo ruido y de las fuertes derivas de línea base añadidas, el modelado químico no colineal preserva características de forma estables que la normalización vectorial resalta, permitiendo una identificación inequívoca del origen del lote.
\newpage
